\paragraph{Why the same trigger uses different shares of the available water.} A suction trigger marks a point on the retention curve. In a sand most available water is released at low suctions, so a $-300$\,cm trigger is reached late (much of the water used); in a clay the water is released over a wide range, so the same suction is reached early.
\paragraph{Why drainage differs under the same rain.} Drainage is the water that exceeds what the profile holds and the crop uses; profiles that hold less (sands) and conduct more pass more water below 1\,m.
\paragraph{Why stage-1 evaporation lasts longer in fine soils.} Fine soils hold more water near the surface and conduct it upward over longer distances, so the surface stays wet longer.
\paragraph{Why the leaf water potential is set by a conductivity floor.} The series resistance of the soil grows without bound as the soil dries; the software bounds the relative conductivity from below to keep the calculation finite, and in dry soil that bound, not the soil, then decides the leaf potential.
\paragraph{Why the drainage share of a mulch saving is understated in coarse soils.} The engines split the saving with ratios that do not contain the soil's drainage competition; in coarse soils, where most saved water drains, that split gives too little to drainage.
\paragraph{Why the meta-model ranks well but has an offset.} It is built from ratios of simulated outcomes, which preserve the order of the scenarios, while its absolute level is anchored to one reference run, so an error in the anchor shifts all values.
\paragraph{Why the daily model fails in sand at dry triggers.} In sand dry triggers lie beyond the suction the crop can produce in the soil, so the full simulation seldom irrigates, while the daily bucket, which has no conductivity limit, still reaches the trigger.

